% ==============================================================================
% =================================== Aufgabe 5 =================================
% ==============================================================================

\chapter*{Zusammenfassung}
\addcontentsline{toc}{chapter}{Zusammenfassung}
\label{chap:Zusammenfassung}
\thispagestyle{fancy}
\section*{Zusammenfassung und Ausblick} 
Die vorliegende Arbeit verdeutlicht, dass moderne \acs{ERP}-Systeme als integratives Rückgrat der betrieblichen Informationsverarbeitung fungieren. Es wurde aufgezeigt, dass die Wandlungsfähigkeit eines Systems die entscheidende Determinante für den langfristigen Unternehmenserfolg in dynamischen Märkten darstellt. Durch den Entwurf eines multidimensionalen Star-Schemas wurde zudem demonstriert, wie komplexe Reseller-Datenstrukturen für analytische Business-Intelligence-Abfragen optimiert werden können. \par\par Ein zentrales Ergebnis ist, dass die Leistungsfähigkeit dieser Systeme maßgeblich von der Datenqualität und der konsequenten Stammdatenzentralisierung abhängt. Der dargestellte Data-Mining-Prozess illustriert ergänzend, wie aus unstrukturierten Internettexten wertvolles Wissen für die strategische Steuerung extrahiert werden kann. Abschließend lässt sich festhalten, dass erst die synergetische Verbindung von operativer Prozessintegration und analytischer Datenauswertung eine objektive Messbarkeit und kontinuierliche Optimierung der unternehmensweiten Wertschöpfungskette ermöglicht.

% ==============================================================================
% =================================== Aufgabe 5 =================================
% ==============================================================================